\chapter{Simulated Annealing (\code{sa})}
\label{ch:sa}

The classical baseline: Metropolis dynamics under a rising-$\beta$ schedule.
Fast, simple, and surprisingly strong on smooth landscapes; every other
engine in \qv{} is measured against it.

\section{The algorithm}

Given a Hamiltonian $E(\cdot)$ (Chapter~\ref{ch:math}) and a schedule
$\beta_1<\dots<\beta_K$:

\begin{keyeq}[Simulated annealing]
\begin{enumerate}
\item Draw a uniformly random initial state $s\in\{-1,+1\}^n$ and compute the
local-field cache $f_i=h_i+\sum_{j\ne i}J_{ij}s_j$.
\item For each schedule step $k=1,\dots,K$, repeat
\code{sweeps\_per\_beta} times:
  \begin{enumerate}
  \item[a.] For each spin $i$ in a freshly shuffled order:
  propose flipping $s_i$; compute $\Delta E=-2s_if_i$
  (Eq.~\eqref{eq:deltaE}); accept with probability
  $\min(1,\mathrm{e}^{-\beta_k\Delta E})$
  (Eq.~\eqref{eq:metropolis}); on acceptance update $s_i\to-s_i$ and the
  affected cache entries $f_j\mathrel{+}=-2s_iJ_{ij}$.
  \end{enumerate}
\item Track the best configuration ever visited; return it with its energy
and the per-step energy trace.
\end{enumerate}
\end{keyeq}

The returned optimum is the \emph{best-ever} state, not the final state ---
a late uphill excursion can never lose an already-found solution.

\paragraph{Complexity.} One sweep costs $O(n)$ proposals; each proposal is
$O(1)$, and each \emph{accepted} flip costs $O(n)$ (dense) or $O(\deg i)$
(sparse) to refresh the cache. A full run is
$O\!\left(K\cdot\code{sweeps\_per\_beta}\cdot n\cdot \bar d\,\right)$ with
$\bar d$ the mean accepted-flip degree --- in practice linear in problem size
for sparse graphs.

\section{Schedule design}
\label{sec:sa-schedules}

\subsection{Fixed linear schedule}

The simplest choice ramps $\beta$ linearly:

\begin{pycode}
from qanneal import AnnealSchedule
schedule = AnnealSchedule.linear(beta_start=0.1, beta_end=5.0, steps=60)
\end{pycode}

Guidance: $\beta_{\text{start}}$ should be hot enough that nearly all moves
are accepted ($\beta_{\text{start}}\Delta \lesssim 0.15$), and
$\beta_{\text{end}}$ cold enough that uphill moves are rare
($\beta_{\text{end}}\Delta \gtrsim 3$), where $\Delta$ is the typical
$|\Delta E|$ of a flip. Since $\Delta$ depends on your coupling scale, fixed
values only suit problems of a known scale --- which motivates:

\subsection{Problem-adaptive tuned schedules (recommended)}
\label{sec:tuned-schedules}

\code{auto\_schedule\_sa\_tuned(problem, mode=...)} removes the guesswork:

\begin{enumerate}
\item Sample \code{probes} (default 256) random single-spin flips on random
states and record $|\Delta E|$ for each.
\item Set the characteristic scale $\Delta$ to the \emph{75th percentile} of
the $|\Delta E|$ sample (robust to outliers, biased toward the barrier-sized
moves that matter).
\item Choose target acceptances $(p_{\text{start}},p_{\text{end}})$ from the
\code{mode} preset and invert Eq.~\eqref{eq:beta-from-p}:
$\beta_{\text{start}}=-\ln p_{\text{start}}/\Delta$,\;
$\beta_{\text{end}}=-\ln p_{\text{end}}/\Delta$.
\item Pick the step count from the preset, growing as $\sqrt n$.
\end{enumerate}

\begin{center}
\small
\begin{tabular}{@{}lcccl@{}}
\toprule
\code{mode} & $p_{\text{start}}$ & $p_{\text{end}}$ & steps & use case\\
\midrule
\code{"fast"}     & 0.85 & 0.22 & $30+6\sqrt n$  & prototyping, large sweeps\\
\code{"balanced"} & 0.88 & 0.10 & $60+6\sqrt n$  & default production\\
\code{"accurate"} & 0.90 & 0.04 & $110+6\sqrt n$ & best solution quality\\
\bottomrule
\end{tabular}
\end{center}

\begin{pycode}
from qanneal import auto_schedule_sa_tuned, solve

schedule = auto_schedule_sa_tuned(ising, mode="balanced")
result   = solve(ising, method="sa", schedule=schedule, reads=32, seed=7)
\end{pycode}

The optional \code{beta\_end\_scale} multiplies $\beta_{\text{end}}$
(values $>1$ freeze harder); \code{steps} overrides the preset count.

\section{Independent replicas: \code{ReplicaAnnealer}}
\label{sec:replica-annealer}

Independent restarts are the cheapest variance reduction: run $R$ chains
from different random seeds and keep the best. \code{ReplicaAnnealer} does
this inside one C++ call, parallelised over replicas with OpenMP:

\begin{pycode}
from qanneal import ReplicaAnnealer

ann = ReplicaAnnealer(ising, schedule, replicas=16)
res = ann.run(sweeps_per_beta=50)
res.global_best_energy       # best over all replicas
res.replicas                 # per-replica results
res.average_energy_trace     # mean energy across replicas per step
\end{pycode}

If a single run finds the optimum with probability $p$, then $R$ independent
runs succeed with probability $1-(1-p)^R$ --- restarts convert a mediocre
per-run success rate into a high aggregate one, at perfectly parallel cost.
The high-level \code{solve(..., reads=R)} applies the same principle across
full solver invocations for every method.

\section{Classical parallel tempering}
\label{sec:classical-pt}

Annealing commits to a one-way trajectory in $\beta$. \emph{Parallel
tempering} (replica exchange) instead runs $R$ chains \emph{simultaneously}
at a fixed ladder $\beta_1<\beta_2<\dots<\beta_R$ and lets configurations
migrate between temperatures.

\subsection{The swap rule, derived}

The joint system samples the product measure
$\Pi(s_1,\dots,s_R)\propto\prod_r \mathrm{e}^{-\beta_r E(s_r)}$.
A \emph{swap move} proposes exchanging the configurations at adjacent rungs
$r,r{+}1$. Detailed balance for the joint measure requires acceptance
\begin{equation}
A_{\text{swap}}
=\min\!\left(1,\;
\frac{\mathrm{e}^{-\beta_r E(s_{r+1})}\,\mathrm{e}^{-\beta_{r+1}E(s_r)}}
     {\mathrm{e}^{-\beta_r E(s_r)}\,\mathrm{e}^{-\beta_{r+1}E(s_{r+1})}}
\right)
\end{equation}
which simplifies to the celebrated form

\begin{keyeq}[Parallel-tempering swap acceptance]
\begin{equation}
A_{\text{swap}}=\min\!\Bigl(1,\;
\mathrm{e}^{\,(\beta_{r+1}-\beta_r)\,\bigl(E(s_{r+1})-E(s_r)\bigr)}\Bigr).
\label{eq:pt-swap}
\end{equation}
\end{keyeq}

A configuration that found a low energy while ``hot'' (small $\beta$, high
mobility) is passed down the ladder toward large $\beta$, where it is
refined; conversely, stuck cold configurations are recycled upward to remelt.
Each chain still satisfies detailed balance at its own $\beta_r$, so all
equilibrium properties are preserved.

\subsection{Usage}

\begin{pycode}
import numpy as np
from qanneal import ParallelTemperingAnnealer

betas = np.linspace(0.1, 5.0, 8).tolist()
ann   = ParallelTemperingAnnealer(ising, betas)
res   = ann.run(sweeps_per_step=50, steps=100, swap_interval=1)

res.best_energy
res.swap_acceptance_trace   # fraction of accepted swaps per step
\end{pycode}

\begin{notebox}[Healthy swap acceptance]
Aim for a mean swap acceptance of roughly $0.2$--$0.5$. Near zero means the
rungs are too far apart (energy distributions at adjacent $\beta$ barely
overlap --- add rungs or shrink the range); near one means rungs are wastefully
close. A geometric $\beta$ ladder is the standard starting point because the
overlap between adjacent rungs is then roughly constant across the ladder.
\end{notebox}

The quantum generalisation of this machinery --- replica exchange on a joint
$(\beta,\Gamma)$ ladder --- is the \code{sqapt} engine of
Chapter~\ref{ch:sqapt}.

\section{When to choose \code{sa}}

\begin{itemize}
\item You need a fast, dependency-free baseline or a sanity check for an
encoding.
\item The landscape is smooth or funnel-like (many greedy paths reach the
optimum).
\item Very large sparse instances where the cost per sweep dominates and the
quantum engines' extra dimension ($M$ slices) is unaffordable.
\end{itemize}

For rugged, frustrated landscapes with tall narrow barriers, continue to
Chapter~\ref{ch:sqa}.
